\section{Software implementation}\label[appendix]{app:code-structure}
The supervised learning implementation for target adaptations is divided into two key components: the nominal control of the robot, programmed in C++ and executed by the Arduino board, and the target adaptation, implemented in Python and managed by an external computer.
The adapter is implemented and trained using Keras with the PyTorch backend. % TODO: cite python and packages?

\subsection{Nominal Control on the Arduino}
Upon startup, the Arduino initializes a connection with the external computer and directs the servo to move to its minimum position.
The Arduino allows some time for the servo to reach this home position and then reads a first measurement from the angle sensor, setting this as the reference $\theta=0$.
During the control loop, the servo reads $\theta$ (relative to the reference) and $\dot{\theta}$ received from the sensor, and reads the target $\theta_d$ received from the external computer.
Based on these measurements, the PID controller computes the value to be written to the servo, but before writing it to the servo, a ceiling is imposed on it to protect the servo from exceeding its mechanical limit (\Cref{subsec:mechanical-details}).
The Arduino allows some time for the servo to act and then reads the measurements from the angle sensor again.
Finally, it sends the tuple containing the state, action, resulting state, and target to the external computer.

\subsection{Target Adaptation on the External Computer}
The target adaptation and supervised learning are implemented exclusively on the external computer through the established connection with the Arduino.
The program is divided into several modules or \emph{threads} that run simultaneously.
This way, values that need to be sent or read can be updated in the background upon request, without delaying the operations of other modules (see~\Cref{app:code-structure} for a more detailed breakdown of the interactions between various threads).
At regular predefined intervals, a \texttt{change\_value} thread changes the true target $\theta_d$ to some value.
Every iteration, it constructs an input feature [$\theta,\,\dot{\theta},\,\theta_d,\,0$] for the adapter.
The resulting $\Delta$ is added to the target, which is read by the \texttt{send} thread and forwarded to the Arduino.
The buffer is maintained by storing the tuples received from the Arduino.
Upon train-time, the adapter weights are updated in the background, without interrupting the adapter's performance.

\Cref{fig:code-diagram} provides an overview of the communication between computer and Arduino, and their inner interactions between various modules.
For example, the sole task of the \texttt{send} thread is to read the value of a global variable representing the target state and send this to the Arduino.
The \texttt{listen} thread reads the tuple sent in return by the Arduino and assigns its contents to the respective shared variables for the state, control action, resulting state and target.
It then adds this tuple to the buffer used for updating the adapter model.
The \texttt{update} thread is responsible for (re)training the adapter.
It creates a copy of the adapter, which can then be trained in the background while the original adapter remains functional.
Once training is completed, the copied adapter's weights are assigned to the original adapter.
The \texttt{change\_value} thread is responsible for setting the ``true'' target in the control loop, as well as the actual adaptations to it.
All threads aim to operate at 50\si{\hertz}.

\begin{figure}[htb]
    \centering
    \def\svgwidth{\linewidth}
    \footnotesize\selectfont\input{Thesis/figures/appendix/code_diagram.pdf_tex}
    \caption{The interaction between the modules (rectangles) and shared variables (ovals) of the computer (blue) and Arduino (red).}
    \label{fig:code-diagram}
\end{figure}



\section{Comparing KOSs}\label[appendix]{app:KOS-refs}

\section{Extended Test Results}\label[appendix]{app:results-ext}